% ID: SEL-D0005-OB-E03
% ESTADO: validada. S006, 2026-11-07; tres pasadas del mismo autor y revisión matemática, didáctica y visual.
\textbf{a) Probabilidad total.}\par\nopagebreak[4]
Sea $H$ ``necesitar hospitalización''. Las causas son excluyentes y exhaustivas; forman una partición. Las tasas de hospitalización son probabilidades condicionadas:
\[
\begin{aligned}
P(H)&=P(A)P(H\mid A)+P(B)P(H\mid B)+P(C)P(H\mid C)\\
&=0.35(0.15)+0.40(0.45)+0.25(0.20)\\
&=0.0525+0.18+0.05=\boxed{0.2825}.
\end{aligned}
\]

\textbf{b) Causa C sabiendo que no necesita hospitalización.}\par\nopagebreak[4]
La condición es $H^c$, no $H$. Tenemos $P(H^c)=1-0.2825=0.7175$ y $P(C\cap H^c)=0.25(0.80)=0.20$. Por Bayes,
\[
P(C\mid H^c)=\frac{0.20}{0.7175}
=\boxed{\frac{80}{287}\approx0.278746}.
\]
Comprobación: las probabilidades conjuntas sin hospitalización son $0.2975,0.22,0.20$, cuya suma es 0.7175. Todas las probabilidades conjuntas de las dos condiciones suman 1. No se puede invertir el condicionamiento intercambiando sus símbolos.
